\documentclass{beamer}
\usepackage{pgfpages}
\setbeameroption{show notes}
\setbeameroption{show notes on second screen=right}

\definecolor{links}{HTML}{2A1B81}
\hypersetup{colorlinks,linkcolor=,urlcolor=links}

\mode<presentation> {

% The Beamer class comes with a number of default slide themes
% which change the colors and layouts of slides. Below this is a list
% of all the themes, uncomment each in turn to see what they look like.

%\usetheme{default}
%\usetheme{AnnArbor}
%\usetheme{Antibes}
%\usetheme{Bergen}
%\usetheme{Berkeley}
%\usetheme{Berlin}
%\usetheme{Boadilla}
\usetheme{CambridgeUS}
%\usetheme{Copenhagen}
%\usetheme{Darmstadt}
%\usetheme{Dresden}
%\usetheme{Frankfurt}
%\usetheme{Goettingen}
%\usetheme{Hannover}
%\usetheme{Ilmenau}
%\usetheme{JuanLesPins}
%\usetheme{Luebeck}
%\usetheme{Madrid}
%\usetheme{Malmoe}
%\usetheme{Marburg}
%\usetheme{Montpellier}
%\usetheme{PaloAlto}
%\usetheme{Pittsburgh}
%\usetheme{Rochester}
%\usetheme{Singapore}
%\usetheme{Szeged}
%\usetheme{Warsaw}

% As well as themes, the Beamer class has a number of color themes
% for any slide theme. Uncomment each of these in turn to see how it
% changes the colors of your current slide theme.

%\usecolortheme{albatross}
%\usecolortheme{beaver}
%\usecolortheme{beetle}
%\usecolortheme{crane}
%\usecolortheme{dolphin}
%\usecolortheme{dove}
%\usecolortheme{fly}
%\usecolortheme{lily}
%\usecolortheme{orchid}
%\usecolortheme{rose}
%\usecolortheme{seagull}
\usecolortheme{seahorse}
%\usecolortheme{whale}
%\usecolortheme{wolverine}

%\setbeamertemplate{footline} % To remove the footer line in all slides uncomment this line
%\setbeamertemplate{footline}[page number] % To replace the footer line in all slides with a simple slide count uncomment this line

%\setbeamertemplate{navigation symbols}{} % To remove the navigation symbols from the bottom of all slides uncomment this line
}
%
%-------------------------------------------------------
% INCLUDE PACKAGES
%-------------------------------------------------------
\usepackage[utf8]{inputenc}
\usepackage[english]{babel}
\usepackage[T1]{fontenc}
\usepackage{helvet}
\usepackage{graphicx} % Allows including images
\usepackage{booktabs} % Allows the use of \toprule, \midrule and \bottomrule in tables
\usepackage{siunitx}
\newcommand{\gc}{\degreeCelsius}
\usepackage{tikz}
%-------------------------------------------------------
% INCLUDE NOTES
%-------------------------------------------------------
%\setbeameroption{show notes}


%-------------------------------------------------------
% INFORMATION IN THE TITLE PAGE
%-------------------------------------------------------

\title[IntSciProg]{\textbf{Tutorial of GNU/Linux in the console. Module IV}}

\subtitle{\textbf{Introduction to Scientific Programming}}


\newcommand\RBox[1]{%
  \tikz\node[draw,rounded corners,align=center,] {#1};%
}  
\author[Curro Pérez-Bernal]
{%
  \texorpdfstring{
    % \begin{columns}
    %   \column{.45\linewidth}
      \centering
      \RBox{Curro Pérez-Bernal\\
        \href{mailto:curropb@uhu.es}{curropb@uhu.es}}
    % \end{columns}
    \vspace{0.5cm}
    \begin{columns}
      \column{.75\linewidth}
      \centering
      \RBox{Integrated Sciences Department --
        University of Huelva}
    \end{columns}
  }
  {John Doe \& Jane Doe}
}
% \author[Jos\'e E. Garc\'{\i}a Ramos]{Jos\'e E. Garc\'{\i}a Ramos}% Your name
% \institute[DCI-UHU]
  
%   \medskip
%   \textit{curropb@uhu.es} % Your email address
  
% \author[Curro P\'erez-Bernal ]{Curro P\'erez Bernal}% Your name
% \institute[DCI-UHU]
  
%   \medskip
%   \textit{curropb@uhu.es} % Your email address

% there must be an empty line above this line - otherwise some unwanted space is added between the university and the country (I do not know why;( )

\date{\today} % Date, can be changed to a custom date

\titlegraphic{\begin{center}\vspace{-0.5cm}\hspace{0.5cm}\includegraphics[width=0.75\linewidth]{Figs/uhu_ceafmc_fisum_isp.png}\end{center}}


%-------------------------------------------------------
% THE BODY OF THE PRESENTATION
%-------------------------------------------------------

\begin{document}

%-------------------------------------------------------
% THE TITLEPAGE
%-------------------------------------------------------

%{\1% % this is the name of the PDF file for the background
\begin{frame}[plain,noframenumbering] % the plain option removes the header from the title page, noframenumbering removes the numbering of this frame only
  \titlepage % call the title page information from above
\end{frame}%}


\begin{frame}{Chapter Outline}{}
\tableofcontents
\end{frame}

\section{System Statistics}
%%%%%%%%%%%%%%%%%%%%%%%%%%%%%%%%%%%%%%%%%%%%%%%%%%%%%%%%%%%%%%%%%%%%%%%%%%%% 
%%%%%%%%%%%%%%%%%%%%%%%%%%%%%%%%%%%%%%%%%%%%%%%%%%%%%%%%%%%%%%%%%%%%%%%%%%%% 
%%%%%%%%%%%%%%%%%%%%%%%%%%%%%%%%%%%%%%%%%%%%%%%%%%%%%%%%%%%%%%%%%%%%%%%%%%%% 
%%%%%%%%%%%%%%%%%%%%%%%%%%%%%%%%%%%%%%%%%%%%%%%%%%%%%%%%%%%%%%%%%%%%%%%%%%%% 
\begin{frame}[t,fragile]{System Statistics: Disk Usage}
  % ------------------------------------------------------


  \begin{block}{The \alert{\texttt{du}} and \alert{\texttt{df}} commands}
    {\footnotesize
The \alert{\texttt{du}} and \alert{\texttt{df}} commands give information on disk usage on the system.    }

  {\scriptsize

  \begin{columns}
      \column{.6\textwidth} The bareword \alert{\texttt{du}} commmand displays the
  disk usage of the specified files and
of each subdirectory (of directory arguments). Using the options
  you can know the size of your disk account or of a particular directory.
      \column{.4\textwidth}
        \hspace{-3cm}
%      {\scriptsize
        \begin{lstlisting}
alice@platea:~$ du .bash*
4	.bash_logout
4	.bashrc
alice@platea:~$ du -hs .
10G     .
        \end{lstlisting}
%      }
    \end{columns}
  \begin{columns}
      \column{.3\textwidth}The \alert{\texttt{df}} command displays the
filesystem disk space  usage for all disks or for a particular system device.
      \column{.7\textwidth}
        \hspace{-3cm}
        {\tiny
        \begin{lstlisting}
alice@platea:~$ df -h
Filesystem  Size  Used Avail Use% Mounted on
/dev/sda3    56G   23G   30G  44% /
tmpfs       2.0G     0  2.0G   0% /lib/init/rw
udev        2.0G  284K  2.0G   1% /dev
tmpfs       2.0G  284K  2.0G   1% /dev/shm
/dev/sda6   134G  7.5G  120G   6% /home
/dev/sda4   2.8G   71M  2.6G   3% /tmp
/dev/sda5   2.8G  702M  2.0G  27% /var
/dev/sdb2   202G   15G  178G   8% /media/usbdisk
alice@platea:~$ 
        \end{lstlisting}
      }
    \end{columns}
  }

  \end{block}
  
%%%%%%%%%%%%%%%%%%%%%%%%%%%%%%%%%%%%%%%%%%%%%%%%%%%%%%%%%%%%%%%%%%%%%%%%%%%%%%%%%%
\note{
{\tiny

Notes Module IV
}
}
\end{frame}
%%%%%%%%%%%%%%%%%%%%%%%%%%%%%%%%%%%%%%%%%%%%%%%%%%%%%%%%%%%%%%%%%%%%%%%%%%%% 
%%%%%%%%%%%%%%%%%%%%%%%%%%%%%%%%%%%%%%%%%%%%%%%%%%%%%%%%%%%%%%%%%%%%%%%%%%%%
\begin{frame}[t,fragile]{Info on User Statistics (I)}
  % ------------------------------------------------------


  \begin{block}{The \alert{\texttt{last}} and \alert{\texttt{uptime}} commands}
    {\footnotesize
   The \alert{\texttt{last}} command shows a listing of last logged in users taken from \texttt{/var/log/wtmp}.

   {\tiny
      \begin{lstlisting}
alice@platea:~$ last
alice  pts/11   deckard.uhu. Mon Apr  9 18:10   still logged in   
migue  pts/11   pris.uhu.es  Mon Apr  9 09:47 - 14:20  (04:33)    
migue  pts/11   pris.uhu.es  Mon Apr  9 09:46 - 09:46  (00:00)    
pepe   pts/11   1.red-5-37.d Wed Apr  4 21:39 - 22:02  (00:23)    
adri   pts/11   85.136.191.d Wed Apr  4 08:59 - 11:13  (02:14)    
ramon  pts/14   144.red-48-7 Tue Apr  3 18:19 - 18:53  (00:33)    
web    pts/11   144.red-48-7 Tue Apr  3 18:19 - 18:53  (00:34)    

wtmp begins Tue Apr  3 18:19:07 2018
alice@platea:~$ 
      \end{lstlisting}
 }
   The \alert{\texttt{uptime}} command tells the current time, how long the system
      has been running, how many users are currently logged on, and
      the system load averages for the past 1, 5, and 15 minutes..


   {\tiny
      \begin{lstlisting}
alice@platea:~$ uptime
 18:10:12 up 207 days, 3:24, 8 users, load average: 0.49, 0.53, 0.40
alice@platea:~$
      \end{lstlisting}
    } 
}
  \end{block}
  
%%%%%%%%%%%%%%%%%%%%%%%%%%%%%%%%%%%%%%%%%%%%%%%%%%%%%%%%%%%%%%%%%%%%%%%%%%%%%%%%%%
\note{
{\tiny

Notes Module IV
}
}
\end{frame}
%%%%%%%%%%%%%%%%%%%%%%%%%%%%%%%%%%%%%%%%%%%%%%%%%%%%%%%%%%%%%%%%%%%%%%%%%%%% 
%%%%%%%%%%%%%%%%%%%%%%%%%%%%%%%%%%%%%%%%%%%%%%%%%%%%%%%%%%%%%%%%%%%%%%%%%%%%
%%%%%%%%%%%%%%%%%%%%%%%%%%%%%%%%%%%%%%%%%%%%%%%%%%%%%%%%%%%%%%%%%%%%%%%%%%%% 
%%%%%%%%%%%%%%%%%%%%%%%%%%%%%%%%%%%%%%%%%%%%%%%%%%%%%%%%%%%%%%%%%%%%%%%%%%%%
\begin{frame}[t,fragile]{Info on User Statistics (II)}
  % ------------------------------------------------------

  \vspace{-0.2cm}
  \begin{block}{The \alert{\texttt{who}} and \alert{\texttt{w}} commands}
    {\footnotesize
   The bareword \alert{\texttt{who}} command shows information about the users currently logged on.
   
   {\tiny
      \begin{lstlisting}
alice@platea:~$ who
alice     tty7     2012-04-18 09:59 (:0)
alice    pts/0     2012-04-18 10:00 (:0.0)
randolph pts/1     2012-04-18 21:06 (atena.us.es)
alice    pts/4     2012-04-18 17:32 (:0.0)
bob      pts/6     2012-04-20 09:01 (platea.uhu.es)
alice@platea:~$
      \end{lstlisting}
 }
   The \alert{\texttt{w}} command      shows information about the users currently logged, the current
    time, the system load, and about what are they currently doing.


   {\tiny
      \begin{lstlisting}
alice@platea:~$ w
 09:17:42 up 1 day, 23:18, 12 users,  load average: 0.00, 0.09, 0.10
USER    TTY     FROM     LOGIN@   IDLE   JCPU   PCPU WHAT
alice   tty7    :0       Wed09   47:18m  9:20m  0.02s xinitrc 
alice   pts/0   :0.0     Wed10   39:46m  8.50s  7.95s ssh tyrell
alice   pts/1   :0.0     Wed13   12:31m  3:53   0.00s less annual.dat
alice   pts/3   :0.0     Wed17   39:45m 42.75s 42.36s ssh rick
alice   pts/5   :0.0     09:01    0.00s  9.87s  0.00s w
alice   pts/9   :0.0     Thu16   16:26m 36.14s  5:34  xfce4-terminal
alice   pts/11  :0.0     Thu19   14:00m  2:05   2:04  soffice.bin
webgem  pts/8   kard.uhu 09:13    3:45   0.01s  0.01s -bash
alice@platea:~$ 
      \end{lstlisting}
    } 
}
  \end{block}
  
%%%%%%%%%%%%%%%%%%%%%%%%%%%%%%%%%%%%%%%%%%%%%%%%%%%%%%%%%%%%%%%%%%%%%%%%%%%%%%%%%%
\note{
{\tiny

Notes Module IV
}
}
\end{frame}
%%%%%%%%%%%%%%%%%%%%%%%%%%%%%%%%%%%%%%%%%%%%%%%%%%%%%%%%%%%%%%%%%%%%%%%%%%%% 
%%%%%%%%%%%%%%%%%%%%%%%%%%%%%%%%%%%%%%%%%%%%%%%%%%%%%%%%%%%%%%%%%%%%%%%%%%%%
%%%%%%%%%%%%%%%%%%%%%%%%%%%%%%%%%%%%%%%%%%%%%%%%%%%%%%%%%%%%%%%%%%%%%%%%%%%% 
%%%%%%%%%%%%%%%%%%%%%%%%%%%%%%%%%%%%%%%%%%%%%%%%%%%%%%%%%%%%%%%%%%%%%%%%%%%%
\section{Filesystem Security}
%%%%%%%%%%%%%%%%%%%%%%%%%%%%%%%%%%%%%%%%%%%%%%%%%%%%%%%%%%%%%%%%%%%%%%%%%%%%
%%%%%%%%%%%%%%%%%%%%%%%%%%%%%%%%%%%%%%%%%%%%%%%%%%%%%%%%%%%%%%%%%%%%%%%%%%%%
%%%%%%%%%%%%%%%%%%%%%%%%%%%%%%%%%%%%%%%%%%%%%%%%%%%%%%%%%%%%%%%%%%%%%%%%%%%%
%%%%%%%%%%%%%%%%%%%%%%%%%%%%%%%%%%%%%%%%%%%%%%%%%%%%%%%%%%%%%%%%%%%%%%%%%%%%
\begin{frame}{UNIX File Permissions}
  % ------------------------------------------------------

  \vspace{-0.3cm}
  \begin{block}{Definition of File Permissions}
    {\footnotesize
    The access rights on files can be checked with the command
\alert{\texttt{ls -l}}.}

  \begin{center}
    \includegraphics[angle=0,width=0.75\linewidth]{./Figs/permissions.png}

  \end{center}
{\scriptsize The permissions, or access rights, and are taken as
  three groups of 3 symbols, affecting the \alert{\emph{file owner}}, \alert{\emph{group member}},
  and \alert{\emph{everyone else}}. 

 The symbols have different meanings if they apply
  to files and directories

\begin{description}
\item[r] permission to read file / list files in dir 
\item[w] permission to modify/delete file / remove, rename, or create files in dir 
\item[x] permission to execute file / enter, run a search or execute files in dir 
\end{description}
}

  \end{block}
  
%%%%%%%%%%%%%%%%%%%%%%%%%%%%%%%%%%%%%%%%%%%%%%%%%%%%%%%%%%%%%%%%%%%%%%%%%%%%%%%%%%
\note{
{\tiny

Notes Module IV
}
}
\end{frame}
%%%%%%%%%%%%%%%%%%%%%%%%%%%%%%%%%%%%%%%%%%%%%%%%%%%%%%%%%%%%%%%%%%%%%%%%%%%%
%%%%%%%%%%%%%%%%%%%%%%%%%%%%%%%%%%%%%%%%%%%%%%%%%%%%%%%%%%%%%%%%%%%%%%%%%%%%
\begin{frame}[t,fragile]{Changing File Permissions}{\emph{Thou know'st 'tis common; all that lives must die...}, WS, Hamlet.}
  % ------------------------------------------------------
  \vspace{-0.3cm}
  \begin{block}{The  \alert{\texttt{chmod}} command}
    {\footnotesize
The  \alert{\texttt{chmod}} command makes changes on file permissions,
also called file mode bits. The syntax of the command is

      \begin{lstlisting}
chmod [ugoa...][[+-=][perms...]...] FILE ...
      \end{lstlisting}
 
Usage Examples}
{\scriptsize

  \begin{columns}
      \column{.5\textwidth}Remove write and read permissions to group and other users.
      \column{.5\textwidth}
        \hspace{-3cm}
%      {\scriptsize
        \begin{lstlisting}
chmod og-rw list_1
        \end{lstlisting}
%      }
    \end{columns}
  \begin{columns}
      \column{.5\textwidth}Give execute permission to the file owner.
      \column{.5\textwidth}
        \hspace{-3cm}
%      {\scriptsize
        \begin{lstlisting}
chmod u+x list_2
        \end{lstlisting}
%      }
    \end{columns}
  \begin{columns}
      \column{.5\textwidth}Give execute permission to all users.
      \column{.5\textwidth}
        \hspace{-3cm}
%      {\scriptsize
        \begin{lstlisting}
chmod a+x li_3 li_4
        \end{lstlisting}
%      }
    \end{columns}
  \begin{columns}
      \column{.5\textwidth}Give group users read and write permissions and remove execute permission.
      \column{.5\textwidth}
        \hspace{-3cm}
%      {\scriptsize
        \begin{lstlisting}
chmod g=rw list_4
        \end{lstlisting}
%      }
    \end{columns}
  }
  \end{block}
  
%%%%%%%%%%%%%%%%%%%%%%%%%%%%%%%%%%%%%%%%%%%%%%%%%%%%%%%%%%%%%%%%%%%%%%%%%%%%%%%%%%
\note{
{\tiny

Notes Module IV
}
}
\end{frame}
%%%%%%%%%%%%%%%%%%%%%%%%%%%%%%%%%%%%%%%%%%%%%%%%%%%%%%%%%%%%%%%%%%%%%%%%%%%% 
%%%%%%%%%%%%%%%%%%%%%%%%%%%%%%%%%%%%%%%%%%%%%%%%%%%%%%%%%%%%%%%%%%%%%%%%%%%%
%%%%%%%%%%%%%%%%%%%%%%%%%%%%%%%%%%%%%%%%%%%%%%%%%%%%%%%%%%%%%%%%%%%%%%%%%%%% 
%%%%%%%%%%%%%%%%%%%%%%%%%%%%%%%%%%%%%%%%%%%%%%%%%%%%%%%%%%%%%%%%%%%%%%%%%%%%
\section{Unix \texttt{Bash} Shell Customization}
%%%%%%%%%%%%%%%%%%%%%%%%%%%%%%%%%%%%%%%%%%%%%%%%%%%%%%%%%%%%%%%%%%%%%%%%%%%%
%%%%%%%%%%%%%%%%%%%%%%%%%%%%%%%%%%%%%%%%%%%%%%%%%%%%%%%%%%%%%%%%%%%%%%%%%%%%
%%%%%%%%%%%%%%%%%%%%%%%%%%%%%%%%%%%%%%%%%%%%%%%%%%%%%%%%%%%%%%%%%%%%%%%%%%%%
%%%%%%%%%%%%%%%%%%%%%%%%%%%%%%%%%%%%%%%%%%%%%%%%%%%%%%%%%%%%%%%%%%%%%%%%%%%%
\begin{frame}[t,fragile]{\texttt{Bash} Shell Variables}
  % ------------------------------------------------------
  \begin{block}{Defining and Exporting Variables in \texttt{Bash}}
    {\footnotesize 
Variables are
parameters that you can create and/or update yourself. Variable names or identifiers consist only of letters, digits and underscores; beginning with a letter or an underscore.

To store \texttt{value} in a variable named \texttt{varname }the following syntax is used:

        \begin{lstlisting}
varname=value
        \end{lstlisting}

To access data stored in a variable we proceed to
  parameter expansion using the special character \alert{\texttt{\$}} as
  in the following example.


        \begin{lstlisting}
alice@platea:~$ foo=bar
alice@platea:~$ echo "foo is $foo..."
foo is bar
\end{lstlisting}
%$

The \texttt{Bash} shell provides automatically a series
of variables called \alert{\emph{environment variables}} that affect the user session and
can be customized.}      
  \end{block}
  
%%%%%%%%%%%%%%%%%%%%%%%%%%%%%%%%%%%%%%%%%%%%%%%%%%%%%%%%%%%%%%%%%%%%%%%%%%%%%%%%%%
\note{
{\tiny

Notes Module IV
}
}
\end{frame}
%%%%%%%%%%%%%%%%%%%%%%%%%%%%%%%%%%%%%%%%%%%%%%%%%%%%%%%%%%%%%%%%%%%%%%%%%%%%
%%%%%%%%%%%%%%%%%%%%%%%%%%%%%%%%%%%%%%%%%%%%%%%%%%%%%%%%%%%%%%%%%%%%%%%%%%%%
\begin{frame}[t,fragile]{\texttt{Bash} Environment Variables}
  % ------------------------------------------------------
\vspace{-0.25cm}  \begin{block}{\alert{Global} and \alert{Local} environment variables}
    {\footnotesize
      Environment variables modify important aspects of the interaction
  between the user and the shell. The \alert{\texttt{printenv}} command
  displays them.}

    {\scriptsize
  \alert{Global} variables are available in all shell sessions and \alert{local} variables are only available in the currently executing shell and will be lost once the user exits the session.

Some of the more important environment variables are the following
\begin{description}
\item[HOME] The current user's home directory, the default for the cd built-in. The value of this variable is also used by tilde expansion.
\item[PATH] A colon-separated list of directories in which the shell looks for commands.
\item[PS1] The primary prompt string.
\item[HISTFILE] The name of the file to which the command history is
  saved. The default value is \texttt{\textasciitilde/.bash\_history}.
\item[HISTFILESIZE] The maximum number of lines contained in the history file, defaults to 500.
\end{description}
}
{\tiny
  \begin{columns}
      \column{.44\textwidth}The redefinition of one of this variables is done as in
  the following example
      \column{.56\textwidth}
        \hspace{-3cm}
%      {\scriptsize
        \begin{lstlisting}
alice@platea:~$ echo $HISTSIZE
500
alice@platea:~$ export HISTSIZE=1000
alice@platea:~$ echo $HISTSIZE
1000
       \end{lstlisting}%$
    \end{columns}
  }
  \end{block}
  
%%%%%%%%%%%%%%%%%%%%%%%%%%%%%%%%%%%%%%%%%%%%%%%%%%%%%%%%%%%%%%%%%%%%%%%%%%%%%%%%%%
\note{
{\tiny
  Both local and global variables are case sensitive and usually capitalised. According to a convention, local variables should be kept lowercase and global variables should be kept uppercase. Although this is just a convention and you are free to keep any name for both local and global variables. A variable name cannot start with a number and should only contain characters in the beginning.

Source:  https://linuxhint.com/bash-environment-variables/
  
Notes Module IV
}
}
\end{frame}
%%%%%%%%%%%%%%%%%%%%%%%%%%%%%%%%%%%%%%%%%%%%%%%%%%%%%%%%%%%%%%%%%%%%%%%%%%%% 
%%%%%%%%%%%%%%%%%%%%%%%%%%%%%%%%%%%%%%%%%%%%%%%%%%%%%%%%%%%%%%%%%%%%%%%%%%%%
\begin{frame}[t,fragile]{\texttt{Bash} User Configuration Files}
  % ------------------------------------------------------

  \vspace{-0.3cm}
  \begin{block}{The \alert{\texttt{\textasciitilde/.bashrc}} config file}
    {\footnotesize
System and user files that give initial values to enviroment vars and can be used to tune
  the shell behavior. The most used file nowadays is \alert{\texttt{\textasciitilde/.bashrc}}. An
  example is shown below.
}

{\scriptsize Other possible configuration files are
  \alert{\texttt{\textasciitilde/.bash\_login}},
  \alert{\texttt{\textasciitilde/.bash\_profile}}, and  \alert{\texttt{\textasciitilde/.profile}}.
}


\vspace{-0.1cm}
{\tiny
        \begin{lstlisting}
# ~/.bashrc: executed by bash(1) for non-login shells.
# see /usr/share/doc/bash/examples/startup-files (in the package bash-doc) for examples

# If not running interactively, don't do anything
case $- in
    *i*) ;;
      *) return;;
esac

# don't put duplicate lines or lines starting with space in the history.
# See bash(1) for more options
HISTCONTROL=ignoreboth
# append to the history file, don't overwrite it
shopt -s histappend
# for setting history length see HISTSIZE and HISTFILESIZE in bash(1)
HISTSIZE=1000
HISTFILESIZE=2000
        \end{lstlisting}%$
}
  \end{block}
  
%%%%%%%%%%%%%%%%%%%%%%%%%%%%%%%%%%%%%%%%%%%%%%%%%%%%%%%%%%%%%%%%%%%%%%%%%%%%%%%%%%
\note{
{\tiny

Notes Module IV
}
}
\end{frame}
%%%%%%%%%%%%%%%%%%%%%%%%%%%%%%%%%%%%%%%%%%%%%%%%%%%%%%%%%%%%%%%%%%%%%%%%%%%%
%%%%%%%%%%%%%%%%%%%%%%%%%%%%%%%%%%%%%%%%%%%%%%%%%%%%%%%%%%%%%%%%%%%%%%%%%%%%
%%%%%%%%%%%%%%%%%%%%%%%%%%%%%%%%%%%%%%%%%%%%%%%%%%%%%%%%%%%%%%%%%%%%%%%%%%%% 
%%%%%%%%%%%%%%%%%%%%%%%%%%%%%%%%%%%%%%%%%%%%%%%%%%%%%%%%%%%%%%%%%%%%%%%%%%%%
\section{Handy Commands}
%%%%%%%%%%%%%%%%%%%%%%%%%%%%%%%%%%%%%%%%%%%%%%%%%%%%%%%%%%%%%%%%%%%%%%%%%%%%
%%%%%%%%%%%%%%%%%%%%%%%%%%%%%%%%%%%%%%%%%%%%%%%%%%%%%%%%%%%%%%%%%%%%%%%%%%%%
%%%%%%%%%%%%%%%%%%%%%%%%%%%%%%%%%%%%%%%%%%%%%%%%%%%%%%%%%%%%%%%%%%%%%%%%%%%%
%%%%%%%%%%%%%%%%%%%%%%%%%%%%%%%%%%%%%%%%%%%%%%%%%%%%%%%%%%%%%%%%%%%%%%%%%%%%
\begin{frame}[t,fragile]{Looking for Files}
  % ------------------------------------------------------

  \vspace{-0.3cm}
  \begin{block}{Using the \alert{\texttt{find}} command}
    {\footnotesize
The \alert{\texttt{find}} command allows to search for files according to different criteria. The syntax of the command is


        \begin{lstlisting}
find [-H] [-L] [-P] [-Olevel] [path] [expression]
        \end{lstlisting}

  Usage Examples
}


\vspace{0.1cm}
{\scriptsize
  \begin{columns}
      \column{.4\textwidth}Find all files with \texttt{txt} suffix in the current directory and its subdirectories.
      \column{.6\textwidth}
        \begin{lstlisting}
find ./ -name "*.txt"
        \end{lstlisting}
    \end{columns}
 \begin{columns}
      \column{.4\textwidth}Find files whose name first character is "c" or "C".
      \column{.6\textwidth}
        \begin{lstlisting}
find ./ -iname "c*"
        \end{lstlisting}
    \end{columns}
 \begin{columns}
      \column{.4\textwidth}Find files in the current folder that are newer than \texttt{filename}.
      \column{.6\textwidth}
        \begin{lstlisting}
find ./ -maxdepth 1 -newer filename
        \end{lstlisting}
    \end{columns}
 \begin{columns}
      \column{.4\textwidth}Find bob user files that are under his home directory and with  group and user write permissions.
      \column{.6\textwidth}
        \begin{lstlisting}
find ~ -user bob -maxdepth 2 \
-perm -u=w,g=w
        \end{lstlisting}
    \end{columns}
}  
  \end{block}
%%%%%%%%%%%%%%%%%%%%%%%%%%%%%%%%%%%%%%%%%%%%%%%%%%%%%%%%%%%%%%%%%%%%%%%%%%%%%%%%%%
\note{
{\tiny

Notes Module IV
}
}
\end{frame}
%%%%%%%%%%%%%%%%%%%%%%%%%%%%%%%%%%%%%%%%%%%%%%%%%%%%%%%%%%%%%%%%%%%%%%%%%%%%
%%%%%%%%%%%%%%%%%%%%%%%%%%%%%%%%%%%%%%%%%%%%%%%%%%%%%%%%%%%%%%%%%%%%%%%%%%%%
\begin{frame}[t,fragile]{Looking into File Contents}
  % ------------------------------------------------------

  \vspace{-0.3cm}
  \begin{block}{Using the \alert{\texttt{grep}} command}
    {\footnotesize
The \alert{\texttt{grep}} command allows to search for files with lines matching a
given pattern.  The syntax of the command is


        \begin{lstlisting}
grep [OPTIONS] pattern [file ... ]
        \end{lstlisting}

  Usage Examples
}


\vspace{0.1cm}
{\scriptsize
  \begin{columns}
      \column{.4\textwidth}Print lines in file \texttt{/etc/group} including the string \texttt{alice}.
      \column{.6\textwidth}
        \begin{lstlisting}
grep alice /etc/group
        \end{lstlisting}
    \end{columns}
 \begin{columns}
      \column{.4\textwidth}Print lines and line numbers in file \texttt{/etc/group} including the string \texttt{alice}.
      \column{.6\textwidth}
        \begin{lstlisting}
grep -n alice /etc/group
        \end{lstlisting}
    \end{columns}
 \begin{columns}
      \column{.4\textwidth}Print lines and line numbers in file \texttt{/etc/group} starting with the string \texttt{alice}.
      \column{.6\textwidth}
        \begin{lstlisting}
grep -c  ^alice /etc/group
        \end{lstlisting}
    \end{columns}
 \begin{columns}
      \column{.4\textwidth}Print number of lines in file \texttt{/etc/group} ending with the string \texttt{alice}.
      \column{.6\textwidth}
        \begin{lstlisting}
grep -c alice$ /etc/group
        \end{lstlisting} %$
    \end{columns}
}  
  \end{block}
%%%%%%%%%%%%%%%%%%%%%%%%%%%%%%%%%%%%%%%%%%%%%%%%%%%%%%%%%%%%%%%%%%%%%%%%%%%%%%%%%%
\note{
{\tiny

Notes Module IV
}
}
\end{frame}
%%%%%%%%%%%%%%%%%%%%%%%%%%%%%%%%%%%%%%%%%%%%%%%%%%%%%%%%%%%%%%%%%%%%%%%%%%%%
%%%%%%%%%%%%%%%%%%%%%%%%%%%%%%%%%%%%%%%%%%%%%%%%%%%%%%%%%%%%%%%%%%%%%%%%%%%%
\begin{frame}[t,fragile]{Comparing File Contents}
  % ------------------------------------------------------

  \vspace{-0.3cm}
  \begin{block}{Using the \alert{\texttt{diff}} command}
    {\footnotesize
Frequently two human-readable files have to be compared to check the differences between them. This can be achieved with the command \alert{\texttt{diff}} that outputs the changes needed to bring one file to the other. The syntax of the command is


        \begin{lstlisting}
diff [OPTION]... FILES
        \end{lstlisting}

  Usage Examples
}


\vspace{0.1cm}
{\scriptsize
  Let's assume that user alice has edited her \texttt{.bashrc} file and want to check the differences with the original copy, prudently saved in file \texttt{mybashrc.bak}.
        \begin{lstlisting}
alice@platea:~$  diff .bashrc mybashrc.bak 
5,6d4
< # If not running interactively, don't do anything
< [ -z "$PS1" ] && return
20c18
< # update the values of LINES and COLUMNS.
---
> # update the values of LINES and COLUMNS. Test.
91,92d88
< #alias la='ls -A'
        \end{lstlisting}
}
  \end{block}
%%%%%%%%%%%%%%%%%%%%%%%%%%%%%%%%%%%%%%%%%%%%%%%%%%%%%%%%%%%%%%%%%%%%%%%%%%%%%%%%%%
\note{
{\tiny

Notes Module IV
}
}
\end{frame}


%%%%%%%%%
\end{document}
